% Options for packages loaded elsewhere
% Options for packages loaded elsewhere
\PassOptionsToPackage{unicode}{hyperref}
\PassOptionsToPackage{hyphens}{url}
%
\documentclass[
  11pt,
  ignorenonframetext,
]{beamer}
\newif\ifbibliography
\usepackage{pgfpages}
\setbeamertemplate{caption}[numbered]
\setbeamertemplate{caption label separator}{: }
\setbeamercolor{caption name}{fg=normal text.fg}
\beamertemplatenavigationsymbolsempty
% remove section numbering
\setbeamertemplate{part page}{
  \centering
  \begin{beamercolorbox}[sep=16pt,center]{part title}
    \usebeamerfont{part title}\insertpart\par
  \end{beamercolorbox}
}
\setbeamertemplate{section page}{
  \centering
  \begin{beamercolorbox}[sep=12pt,center]{section title}
    \usebeamerfont{section title}\insertsection\par
  \end{beamercolorbox}
}
\setbeamertemplate{subsection page}{
  \centering
  \begin{beamercolorbox}[sep=8pt,center]{subsection title}
    \usebeamerfont{subsection title}\insertsubsection\par
  \end{beamercolorbox}
}
% Prevent slide breaks in the middle of a paragraph
\widowpenalties 1 10000
\raggedbottom
\AtBeginPart{
  \frame{\partpage}
}
\AtBeginSection{
  \ifbibliography
  \else
    \frame{\sectionpage}
  \fi
}
\AtBeginSubsection{
  \frame{\subsectionpage}
}
\usepackage{iftex}
\ifPDFTeX
  \usepackage[T1]{fontenc}
  \usepackage[utf8]{inputenc}
  \usepackage{textcomp} % provide euro and other symbols
\else % if luatex or xetex
  \usepackage{unicode-math} % this also loads fontspec
  \defaultfontfeatures{Scale=MatchLowercase}
  \defaultfontfeatures[\rmfamily]{Ligatures=TeX,Scale=1}
\fi
\usepackage{lmodern}

\usetheme[]{Madrid}
\ifPDFTeX\else
  % xetex/luatex font selection
\fi
% Use upquote if available, for straight quotes in verbatim environments
\IfFileExists{upquote.sty}{\usepackage{upquote}}{}
\IfFileExists{microtype.sty}{% use microtype if available
  \usepackage[]{microtype}
  \UseMicrotypeSet[protrusion]{basicmath} % disable protrusion for tt fonts
}{}
\makeatletter
\@ifundefined{KOMAClassName}{% if non-KOMA class
  \IfFileExists{parskip.sty}{%
    \usepackage{parskip}
  }{% else
    \setlength{\parindent}{0pt}
    \setlength{\parskip}{6pt plus 2pt minus 1pt}}
}{% if KOMA class
  \KOMAoptions{parskip=half}}
\makeatother


\usepackage{longtable,booktabs,array}
\usepackage{calc} % for calculating minipage widths
\usepackage{caption}
% Make caption package work with longtable
\makeatletter
\def\fnum@table{\tablename~\thetable}
\makeatother
\usepackage{graphicx}
\makeatletter
\newsavebox\pandoc@box
\newcommand*\pandocbounded[1]{% scales image to fit in text height/width
  \sbox\pandoc@box{#1}%
  \Gscale@div\@tempa{\textheight}{\dimexpr\ht\pandoc@box+\dp\pandoc@box\relax}%
  \Gscale@div\@tempb{\linewidth}{\wd\pandoc@box}%
  \ifdim\@tempb\p@<\@tempa\p@\let\@tempa\@tempb\fi% select the smaller of both
  \ifdim\@tempa\p@<\p@\scalebox{\@tempa}{\usebox\pandoc@box}%
  \else\usebox{\pandoc@box}%
  \fi%
}
% Set default figure placement to htbp
\def\fps@figure{htbp}
\makeatother





\setlength{\emergencystretch}{3em} % prevent overfull lines

\providecommand{\tightlist}{%
  \setlength{\itemsep}{0pt}\setlength{\parskip}{0pt}}



 


\makeatletter
\@ifpackageloaded{caption}{}{\usepackage{caption}}
\AtBeginDocument{%
\ifdefined\contentsname
  \renewcommand*\contentsname{Table of contents}
\else
  \newcommand\contentsname{Table of contents}
\fi
\ifdefined\listfigurename
  \renewcommand*\listfigurename{List of Figures}
\else
  \newcommand\listfigurename{List of Figures}
\fi
\ifdefined\listtablename
  \renewcommand*\listtablename{List of Tables}
\else
  \newcommand\listtablename{List of Tables}
\fi
\ifdefined\figurename
  \renewcommand*\figurename{Figure}
\else
  \newcommand\figurename{Figure}
\fi
\ifdefined\tablename
  \renewcommand*\tablename{Table}
\else
  \newcommand\tablename{Table}
\fi
}
\@ifpackageloaded{float}{}{\usepackage{float}}
\floatstyle{ruled}
\@ifundefined{c@chapter}{\newfloat{codelisting}{h}{lop}}{\newfloat{codelisting}{h}{lop}[chapter]}
\floatname{codelisting}{Listing}
\newcommand*\listoflistings{\listof{codelisting}{List of Listings}}
\makeatother
\makeatletter
\makeatother
\makeatletter
\@ifpackageloaded{caption}{}{\usepackage{caption}}
\@ifpackageloaded{subcaption}{}{\usepackage{subcaption}}
\makeatother

\usepackage{bookmark}
\IfFileExists{xurl.sty}{\usepackage{xurl}}{} % add URL line breaks if available
\urlstyle{same}
\hypersetup{
  pdftitle={How LLMs Form Input \& Output Tensors},
  hidelinks,
  pdfcreator={LaTeX via pandoc}}


\title{How LLMs Form Input \& Output Tensors}
\subtitle{Context Length • Token Windows • Embeddings • Batch Size}
\author{}
\date{}

\begin{document}
\frame{\titlepage}


\begin{frame}{Overview}
\phantomsection\label{overview}
Concepts covered:

\begin{itemize}
\tightlist
\item
  Tokenization of text\\
\item
  Context length and token windows\\
\item
  Input--output structure for next-token prediction\\
\item
  Embedding dimension and matrix form\\
\item
  Batch size and tensor shapes
\end{itemize}
\end{frame}

\begin{frame}{Tokenized Example Sentence}
\phantomsection\label{tokenized-example-sentence}
We use the sentence:

``I love large language models.''

Token IDs (GPT-4o tokenizer):

{[} {[}40,; 584,; 2201,; 7457,; 9937,; 13{]} {]}

Define the sequence:

{[} t\_1 = 40,\quad t\_2 = 584,\quad t\_3 = 2201,\quad t\_4 = 7457,\quad
t\_5 = 9937,\quad t\_6 = 13. {]}
\end{frame}

\begin{frame}{Context Length = 2}
\phantomsection\label{context-length-2}
The model uses a \textbf{context window of size (2)}.\\
Training pairs for next-token prediction are:

{[}

\begin{aligned}
(t_1,\, t_2) &\rightarrow t_3, \\
(t_2,\, t_3) &\rightarrow t_4, \\
(t_3,\, t_4) &\rightarrow t_5, \\
(t_4,\, t_5) &\rightarrow t_6.
\end{aligned}

{]}

This corresponds to supervised learning:

{[} X\_i = (t\_i,, t\_\{i+1\}), \qquad y\_i = t\_\{i+2\}. {]}
\end{frame}

\begin{frame}{Embedding Dimension and Input Matrices}
\phantomsection\label{embedding-dimension-and-input-matrices}
Let the embedding dimension be (d).\\
Each token (t) maps to an embedding vector:

{[} E{[}t{]} \in \mathbb{R}\^{}d. {]}

Therefore each input (X\_i) becomes a matrix in (\mathbb{R}\^{}\{2
\times d\}):

{[} X\_i =

\begin{bmatrix}
E[t_i] \\
E[t_{i+1}]
\end{bmatrix}

, \qquad X\_i \in \mathbb{R}\^{}\{2 \times d\}. {]}
\end{frame}

\begin{frame}{All Four Training Examples}
\phantomsection\label{all-four-training-examples}
{[}

\begin{aligned}
X_1 &= 
\begin{bmatrix}
E[t_1]\\[2pt] E[t_2]
\end{bmatrix},
& y_1 &= t_3,
\\[6pt]
X_2 &= 
\begin{bmatrix}
E[t_2]\\[2pt] E[t_3]
\end{bmatrix},
& y_2 &= t_4,
\\[6pt]
X_3 &= 
\begin{bmatrix}
E[t_3]\\[2pt] E[t_4]
\end{bmatrix},
& y_3 &= t_5,
\\[6pt]
X_4 &= 
\begin{bmatrix}
E[t_4]\\[2pt] E[t_5]
\end{bmatrix},
& y_4 &= t_6.
\end{aligned}

{]}

Each (X\_i) is a (2 \times d) matrix.
\end{frame}

\begin{frame}{Batch Size = 1}
\phantomsection\label{batch-size-1}
If we train one example at a time:

{[} X\_i \in \mathbb{R}\^{}\{2 \times d\}, \qquad y\_i \in \mathbb{N}.
{]}

No combined tensor is formed.\\
Each example is processed independently.
\end{frame}

\begin{frame}{Batch Size = 4 (All Samples Together)}
\phantomsection\label{batch-size-4-all-samples-together}
If we train on all four samples simultaneously:

{[} X =

\begin{bmatrix}
X_1 \\[4pt]
X_2 \\[4pt]
X_3 \\[4pt]
X_4
\end{bmatrix}

\in \mathbb{R}\^{}\{4 \times 2 \times d\}. {]}

The target vector:

{[} y = (t\_3,; t\_4,; t\_5,; t\_6) \in \mathbb{N}\^{}4. {]}
\end{frame}

\begin{frame}{Tensor Shape Summary}
\phantomsection\label{tensor-shape-summary}
Context length: (2)\\
Embedding dimension: (d)\\
Number of samples: (4)

Batch size (=1):

{[} X\_i \in \mathbb{R}\^{}\{2 \times d\}, \qquad y\_i \in \mathbb{N}.
{]}

Batch size (=4):

{[} X \in \mathbb{R}\^{}\{4 \times 2 \times d\}, \qquad y
\in \mathbb{N}\^{}4. {]}
\end{frame}

\begin{frame}{General Case}
\phantomsection\label{general-case}
For context length (k):

{[} X\_i \in \mathbb{R}\^{}\{k \times d\}. {]}

For batch size (B):

{[} X \in \mathbb{R}\^{}\{B \times k \times d\}, \qquad y
\in \mathbb{N}\^{}B. {]}
\end{frame}

\begin{frame}{Key Takeaways}
\phantomsection\label{key-takeaways}
\begin{itemize}
\tightlist
\item
  Tokens form a sequence ((t\_1, \dots, t\_n)).\\
\item
  Context windows create supervised learning pairs.\\
\item
  Embeddings convert tokens to vectors in (\mathbb{R}\^{}d).\\
\item
  Each input is a matrix (X\_i \in \mathbb{R}\^{}\{k \times d\}).\\
\item
  Batch size determines the shape of the full tensor.
\end{itemize}
\end{frame}

\begin{frame}{Thank You!}
\phantomsection\label{thank-you}
I can also add:

\begin{itemize}
\tightlist
\item
  attention equations\\
\item
  positional encoding math\\
\item
  matrix illustrations\\
\item
  model pipeline diagrams
\end{itemize}

Just tell me!
\end{frame}




\end{document}
